\[ %%%%%%%%%%%%%%%%%%%%%%%%%%%%%%%%%%%%%%%%%%%%%%%%%%%%%%%%%%%%%%%%%%%%%%%%%%%%%%%% \subsection{Stability-Regularized Optimization} \label{sec:stability_optimization} %%%%%%%%%%%%%%%%%%%%%%%%%%%%%%%%%%%%%%%%%%%%%%%%%%%%%%%%%%%%%%%%%%%%%%%%%%%%%%%% Long-horizon prediction of nonlinear dynamical systems requires not only accurate state estimation but also stable evolution of the learned dynamical operator. A model optimized exclusively for prediction accuracy may produce physically incorrect trajectories due to accumulated approximation errors during recursive rollout. To address this limitation, PHYSGEN introduces stability-regularized optimization, where the learned evolution operator is explicitly constrained to generate bounded and dynamically consistent trajectories. %%%%%%%%%%%%%%%%%%%%%%%%%%%%%%%%%%%%%%%%%%%%%%%%%%%%%%%%%%%%%%%%%%%%%%%%%%%%%%%% \subsubsection{Stability Objective} %%%%%%%%%%%%%%%%%%%%%%%%%%%%%%%%%%%%%%%%%%%%%%%%%%%%%%%%%%%%%%%%%%%%%%%%%%%%%%%% The standard forecasting objective can be extended with a stability term: \begin{equation} \theta^* = \arg\min_\theta \left( \mathcal{L}_{forecast} + \lambda_{stab} \mathcal{L}_{stab} \right), \end{equation} where: \begin{itemize} \item $\mathcal{L}_{forecast}$ measures prediction accuracy; \item $\mathcal{L}_{stab}$ regulates dynamical stability; \item $\lambda_{stab}$ controls the strength of stability regularization. \end{itemize} %%%%%%%%%%%%%%%%%%%%%%%%%%%%%%%%%%%%%%%%%%%%%%%%%%%%%%%%%%%%%%%%%%%%%%%%%%%%%%%% \subsubsection{Latent State Stability Constraint} %%%%%%%%%%%%%%%%%%%%%%%%%%%%%%%%%%%%%%%%%%%%%%%%%%%%%%%%%%%%%%%%%%%%%%%%%%%%%%%% Let $H_t$ denote the latent physical state generated by PHYSGEN. The uncontrolled evolution is: \begin{equation} H_{t+1} = \Psi_\theta(H_t). \end{equation} For stable dynamics, the latent state variation should remain bounded: \begin{equation} \| H_{t+1}-H_t \| < \delta, \end{equation} where $\delta$ represents an admissible state transition magnitude. The latent stability loss is defined as: \begin{equation} \mathcal{L}_{latent} = \| H_{t+1}-H_t \|^2. \end{equation} This constraint prevents uncontrolled growth of latent representations. %%%%%%%%%%%%%%%%%%%%%%%%%%%%%%%%%%%%%%%%%%%%%%%%%%%%%%%%%%%%%%%%%%%%%%%%%%%%%%%% \subsubsection{Lyapunov-Inspired Regularization} %%%%%%%%%%%%%%%%%%%%%%%%%%%%%%%%%%%%%%%%%%%%%%%%%%%%%%%%%%%%%%%%%%%%%%%%%%%%%%%% The stability formulation of PHYSGEN is inspired by Lyapunov stability theory. A Lyapunov function $V(H)$ satisfies: \begin{equation} V(H_{t+1})-V(H_t)\leq0 \end{equation} for stable system evolution. PHYSGEN does not require an analytical Lyapunov function. Instead, it learns a latent stability energy: \begin{equation} V_\theta(H_t). \end{equation} The corresponding regularization term is: \begin{equation} \mathcal{L}_{Lyap} = \max ( 0, V_\theta(H_{t+1}) - V_\theta(H_t) ) . \end{equation} This formulation penalizes only unstable energy growth. %%%%%%%%%%%%%%%%%%%%%%%%%%%%%%%%%%%%%%%%%%%%%%%%%%%%%%%%%%%%%%%%%%%%%%%%%%%%%%%% \subsubsection{Trajectory Divergence Control} %%%%%%%%%%%%%%%%%%%%%%%%%%%%%%%%%%%%%%%%%%%%%%%%%%%%%%%%%%%%%%%%%%%%%%%%%%%%%%%% For long-horizon forecasting, the predicted trajectory is: \begin{equation} \hat{X}_{t+k} = \Phi_\theta^k(X_t). \end{equation} The divergence between predicted and reference trajectories is: \begin{equation} D_k = \| X_{t+k} - \hat{X}_{t+k} \|. \end{equation} PHYSGEN introduces a trajectory growth penalty: \begin{equation} \mathcal{L}_{div} = \sum_{k=1}^{K} \rho_k D_k, \end{equation} where $\rho_k$ controls the contribution of different prediction horizons. This encourages bounded error propagation during recursive simulation. %%%%%%%%%%%%%%%%%%%%%%%%%%%%%%%%%%%%%%%%%%%%%%%%%%%%%%%%%%%%%%%%%%%%%%%%%%%%%%%% \subsubsection{Adaptive Stability Controller Integration} %%%%%%%%%%%%%%%%%%%%%%%%%%%%%%%%%%%%%%%%%%%%%%%%%%%%%%%%%%%%%%%%%%%%%%%%%%%%%%%% The stability regularization is coupled with the Adaptive Stability Control Module introduced in Section~\ref{sec:stability_control}. The corrected latent evolution becomes: \begin{equation} H_{t+1} = \Psi_\theta(H_t) + C_{stab}(H_t), \end{equation} where the correction term is optimized to minimize instability: \begin{equation} C_{stab} = -\nabla_H \mathcal{L}_{stab}. \end{equation} Therefore, stability is not only a training constraint but an active component of the prediction process. %%%%%%%%%%%%%%%%%%%%%%%%%%%%%%%%%%%%%%%%%%%%%%%%%%%%%%%%%%%%%%%%%%%%%%%%%%%%%%%% \subsubsection{Complete Stability-Regularized Objective} %%%%%%%%%%%%%%%%%%%%%%%%%%%%%%%%%%%%%%%%%%%%%%%%%%%%%%%%%%%%%%%%%%%%%%%%%%%%%%%% The complete stability objective is defined as: \begin{equation} \mathcal{L}_{stab} = \alpha \mathcal{L}_{latent} + \beta \mathcal{L}_{Lyap} + \gamma \mathcal{L}_{div} + \eta \mathcal{L}_{phys}, \end{equation} where: \begin{itemize} \item $\mathcal{L}_{latent}$ controls latent state variation; \item $\mathcal{L}_{Lyap}$ constrains stability energy growth; \item $\mathcal{L}_{div}$ limits trajectory divergence; \item $\mathcal{L}_{phys}$ maintains physical consistency. \end{itemize} The final PHYSGEN optimization problem becomes: \begin{equation} \boxed{ \theta^* = \arg\min_\theta \left[ \mathcal{L}_{data} + \lambda_p\mathcal{L}_{physics} + \lambda_g\mathcal{L}_{graph} + \lambda_r\mathcal{L}_{rollout} + \lambda_s ( \alpha\mathcal{L}_{latent} + \beta\mathcal{L}_{Lyap} + \gamma\mathcal{L}_{div} ) \right] } \end{equation} %%%%%%%%%%%%%%%%%%%%%%%%%%%%%%%%%%%%%%%%%%%%%%%%%%%%%%%%%%%%%%%%%%%%%%%%%%%%%%%% \subsubsection{Stability Interpretation} %%%%%%%%%%%%%%%%%%%%%%%%%%%%%%%%%%%%%%%%%%%%%%%%%%%%%%%%%%%%%%%%%%%%%%%%%%%%%%%% The stability-regularized formulation changes the learning objective from trajectory fitting to controlled dynamical system identification. The learned operator is encouraged to satisfy: \begin{equation} \Phi_\theta: \mathcal{M}_{phys} \rightarrow \mathcal{M}_{phys}, \end{equation} while maintaining bounded long-term evolution: \begin{equation} \lim_{k\rightarrow\infty} \| \hat{X}_{t+k} \| < \infty. \end{equation} %%%%%%%%%%%%%%%%%%%%%%%%%%%%%%%%%%%%%%%%%%%%%%%%%%%%%%%%%%%%%%%%%%%%%%%%%%%%%%%% \subsubsection{Summary} %%%%%%%%%%%%%%%%%%%%%%%%%%%%%%%%%%%%%%%%%%%%%%%%%%%%%%%%%%%%%%%%%%%%%%%%%%%%%%%% The stability-regularized optimization framework provides PHYSGEN with an explicit mechanism for controlling long-horizon prediction behavior. By integrating Lyapunov-inspired constraints, latent-state regulation, and adaptive stability correction, PHYSGEN transforms stability from an implicit property of training into an explicit architectural objective. This formulation is essential for reliable prediction of complex nonlinear systems where small errors can otherwise accumulate into physically invalid trajectories. \] 
